\documentclass[12pt,a4paper]{article}
\usepackage[utf8]{inputenc}
\usepackage[T1]{fontenc}
\usepackage[italian]{babel}
\usepackage{hyperref}
\usepackage{tabularx}
\usepackage{booktabs}
\usepackage{geometry}
\geometry{a4paper, margin=2cm}

\title{Contextual Inquiry: FarFromHome}
\author{}
\date{}

\begin{document}
\maketitle

\section{Introduzione al progetto}
L'app \textbf{AvanziZero} è rivolta a coinquilini, chi vive solo e famiglie, posizionandosi come una soluzione avanzata per la gestione collaborativa della dispensa e l'ottimizzazione degli acquisti (lista della spesa) per minimizzare gli sprechi alimentari, evitare acquisti ridondanti e ridurre il carico cognitivo legato alla pianificazione.

L'app si articola attorno a quattro macro-funzionalità:
\begin{enumerate}
    \item \textbf{Autenticazione \& Gestione Gruppi Casa:} registrazione, creazione di un gruppo con codice univoco, accesso a un gruppo esistente via codice condiviso, sistema di richieste di adesione, ruoli Admin/Membro.
    \item \textbf{Dispensa Digitale:} inventario della dispensa in tempo reale, condiviso tra tutti i membri del gruppo, con scansione degli scontrini, filtri per categoria, barra di ricerca e alert per le scadenze.
    \item \textbf{Lista della Spesa:} lista digitale condivisa tra i coinquilini, filtrabile per categoria, con spunta degli elementi acquistati e trasferimento automatico in dispensa.
    \item Il valore differenziante del progetto risiede nell'integrazione di un motore di \textbf{Intelligenza Artificiale Predittiva}.
\end{enumerate}

Per evitare dimenticanze e sprechi di tempo, sono introdotte funzionalità che analizzino le abitudini di consumo del gruppo inserendo automaticamente gli articoli da inserire nella lista della spesa. L'utente, aprendo l'app, troverà già una lista di suggerimenti intelligenti basati sulle scadenze e sull'esaurimento dei prodotti e, al ritorno, potrà inserire automaticamente la spesa in dispensa scannerizzando lo scontrino (un'altra IA riconoscerà i prodotti acquistati con le relative quantità e li inserirà nella dispensa), riducendo il carico cognitivo e il tempo speso a controllare manualmente la dispensa e inserire i prodotti.

\section{Problem scenario}

\subsection{Scenario 1: Il fuorisede sbadato}
\textbf{Attore:} Marco è uno studente universitario fuorisede che condivide l'appartamento con altri due ragazzi. Ha una vita frenetica tra lezioni, studio e palestra. Spesso fa la spesa all'ultimo minuto mentre torna a casa.

\textbf{Scenario:} Marco è al supermercato e cerca di ricordarsi cosa manca a casa. Apre WhatsApp per cercare i messaggi dei coinquilini, ma la chat è disordinata. Cerca di fare mente locale su cosa ha mangiato in settimana. Fatta la spesa e tornato a casa, si rende conto di aver comprato la pasta (che in realtà c'era già) ma di aver dimenticato l'olio, fondamentale per la cena, e che un altro coinquilino aveva già comprato le uova. Pensa che sarebbe fantastico sapere cosa c'è già in casa e cosa è stato comprato senza dover chiedere ogni volta ai coinquilini.

\vspace{0.5cm}
\noindent\textbf{Claims Scenario 1}

\begin{table}[h]
\centering
\begin{tabularx}{\textwidth}{|l|X|}
\hline
\textbf{Situation Features} & \textbf{Pros(+) e Cons(-)} \\
\hline
Spesa individuale senza lista condivisa, & \textbf{+} Nessuna dipendenza da strumenti digitali. \\
affidata alla memoria & \textbf{-} Impossibilità di coordinarsi in tempo reale a distanza. \\
& \textbf{-} Alto rischio di dimenticare articoli essenziali. \\
& \textbf{-} Alta probabilità di acquisti doppi. \\
\hline
\end{tabularx}
\end{table}

\subsection{Scenario 2: La mamma smemorata}
\textbf{Attore:} Giulia è una mamma lavoratrice. È lei che, in famiglia, gestisce la logistica della casa. Ama avere tutto sotto controllo ma, avendo tante cose per la testa, spesso dimentica qualche scadenza e il cibo che c'è in casa.

\textbf{Scenario:} Giulia deve preparare una torta, prende il latte ma si accorge che è scaduto da due giorni. Decide di cambiare ricetta e prende il mascarpone dal ripiano più vicino, accorgendosi poi di aver usato quello a scadenza più lunga e non il mascarpone che scadeva quel giorno. Inoltre voleva usare la panna, ricordando di averne una. Tuttavia l'aveva usata il giorno prima e, non avendola aggiunta alla lista della spesa, si era dimenticata di comprarla. Giulia vorrebbe un sistema che la aiutasse a ricordare le scadenze del cibo per minimizzare gli sprechi, e che le ricordasse di aggiungere gli elementi alla lista della spesa.

\vspace{0.5cm}
\noindent\textbf{Claims Scenario 2}

\begin{table}[h]
\centering
\begin{tabularx}{\textwidth}{|l|X|}
\hline
\textbf{Situation Features} & \textbf{Pros(+) e Cons(-)} \\
\hline
Ricordare le scadenze e segnare & \textbf{+} Massimo controllo su ciò che si possiede. \\
le cose da comprare nella lista & \textbf{-} Metodo inadatto per chi è molto distratto. \\
della spesa & \textbf{-} Se non si ha la lista sottomano è molto probabile dimenticare di inserire i prodotti. \\
\hline
\end{tabularx}
\end{table}

\subsection{Scenario 3: L'inventario manuale che non si fa mai}
\textbf{Attore:} Elena, 38 anni, project manager molto pigra che si è appena trasferita in una nuova città da sola.

\textbf{Scenario:} Ogni volta che deve fare la spesa Elena controlla fisicamente frigorifero, dispensa e stipetti per capire cosa serve. Impiega 15-20 minuti per questo inventario, dopodiché redige la lista su carta. Impiega altro tempo a controllare di cosa ha meno scorte, in modo da non dover andare più volte al supermercato. Dopo essere tornata deve poi ricordare tutto quello che ha comprato, per essere sicura di consumarlo prima della scadenza, e ogni volta deve reinserire ogni prodotto. Vorrebbe che qualcuno, al posto suo, controllasse cosa sta per finire le preparasse già la lista. Vorrebbe inoltre sapere dove sono i supermercati più vicini, in modo da non dover cercare ogni volta online e non essere costretta ad andare nell'unico che conosce che è abbastanza lontano da casa.

\vspace{0.5cm}
\noindent\textbf{Claims Scenario 3}

\begin{table}[h]
\centering
\begin{tabularx}{\textwidth}{|l|X|}
\hline
\textbf{Situation Features} & \textbf{Pros(+) e Cons(-)} \\
\hline
Lista cartacea manuale e & \textbf{+} Controllo diretto e accurato sull'inventario fisico. \\
viaggio lungo verso il & \textbf{+} Conoscenza del supermercato in cui si va. \\
supermercato & \textbf{-} Processo lento, ripetitivo e ad alto carico cognitivo. \\
& \textbf{-} Non avvisa delle scadenze imminenti e dei prodotti che stanno per finire. \\
& \textbf{-} Viaggio verso il supermercato lungo. \\
& \textbf{-} Possibilità di non conoscere i supermercati in una nuova zona. \\
\hline
\end{tabularx}
\end{table}

\section{Piano di indagine contestuale}
\textit{Definisci chi andrai ad intervistare e con quali obiettivi.}

Al fine di comprendere l'importanza del problema (dimenticare prodotti o perdere tempo a fare la lista) e validare la soluzione IA, l'indagine si rivolge agli stakeholder principali: \textbf{chi fa la spesa per un gruppo o per sé stesso}.
Sono state individuate due tipologie:
\begin{enumerate}
    \item Il coinquilino/fuorisede (spesa disorganizzata/condivisa)
    \item L'organizzatore della famiglia/coppia (spesa pianificata)
\end{enumerate}

\textbf{Gli obiettivi dell'analisi:}
\begin{itemize}
    \item Capire quanto tempo viene speso per compilare la lista della spesa.
    \item Identificare la frequenza con cui ci si dimentica di comprare prodotti essenziali.
    \item Scoprire come vengono attualmente tracciate le abitudini di consumo.
    \item Valutare la propensione degli utenti ad affidarsi ai suggerimenti di un'Intelligenza Artificiale.
\end{itemize}

\subsection{Prima intervista (Screening iniziale)}
Seleziona un piccolo campione eterogeneo (es. 5 persone: 2 fuorisede, 2 conviventi, 1 single lavoratore) per porre domande generali (anche online/form):
\begin{itemize}
    \item Chi fa solitamente la spesa in casa tua?
    \item Quanto tempo impieghi a preparare la lista della spesa?
    \item Ti è mai capitato di tornare dal supermercato e accorgerti di aver dimenticato un prodotto essenziale? Quanto spesso accade?
    \item Come comunichi con i tuoi coinquilini/familiari su cosa comprare?
    \item Ti darebbe fastidio se un'app analizzasse i tuoi consumi alimentari per aiutarti?
\end{itemize}

\subsection{Seconda intervista (Contextual Inquiry vera e propria - Face-to-Face)}
\textit{(Questa è la parte cruciale. Devi osservare gli utenti \textbf{mentre} creano una lista della spesa o fargli domande di contesto).}

\textbf{Presentazione:} "Stiamo sviluppando FarFromHome, un'app che usa l'IA per suggerirti automaticamente cosa manca in dispensa. Ci serve capire le tue abitudini attuali."

\textit{Durante l'osservazione (o intervista diretta), chiedi all'utente di simulare come fa la lista della spesa ORA e fagli domande contestuali:}
\begin{itemize}
    \item "Fammi vedere come hai annotato le cose da comprare per oggi."
    \item "Perché hai aggiunto il latte? Come ti sei accorto che mancava?" \textit{(Per capire se vanno a memoria o guardano nel frigo)}
    \item "Cosa fai se sei al supermercato e non sei sicuro se a casa c'è ancora della passata di pomodoro?"
    \item "Se un'app ti proponesse un tasto 'Aggiungi i soliti prodotti della settimana', lo useresti o vorresti comunque controllare uno per uno?"
\end{itemize}

\subsection{Modalità di somministrazione}
Seguendo il template, specifica che hai somministrato il questionario \textit{face-to-face} con due persone del tuo team (uno fa le domande, l'altro prende appunti e osserva il linguaggio del corpo/reazioni, ad esempio se storcono il naso all'idea dell'IA). Mostra una tabella con i profili degli intervistati (Sesso, Età, Lavoro, Status abitativo es. Fuorisede/Famiglia).

\section{Analisi dei risultati delle indagini contestuali}

\subsection{Quali attività svolgono al momento gli utenti?}
Al momento gli utenti, per capire cosa comprare, aprono fisicamente il frigorifero e gli stipetti (inventario visivo). Dopodiché scrivono i prodotti su chat WhatsApp condivise (es. gruppo "Spesa Casa") o su app di note generiche (es. Google Keep). Quando un prodotto finisce durante la settimana, sperano di ricordarsi di appuntarlo, ma spesso se ne dimenticano.

\subsection{Quali task vorrebbero svolgere?}
Vorrebbero aprire l'app e trovare \textbf{già pronta} una bozza della lista della spesa con i prodotti che stanno per finire, senza dover controllare fisicamente la dispensa. Vogliono poter accettare i suggerimenti dell'IA con un solo tap e avere la libertà di aggiungere manualmente prodotti extra.

\subsection{Come vengono apprese le attività da svolgere?}
L'interazione con l'IA deve essere passiva e \textit{zero-effort}. L'utente non deve "imparare" a usare l'IA: l'app impara dall'utente man mano che vengono inseriti o consumati i prodotti. L'attività di accettare i suggerimenti sarà intuitiva, simile a spuntare le notifiche o accettare suggerimenti testuali sulla tastiera.

\subsection{Dove vengono svolte le attività?}
\begin{itemize}
    \item \textbf{Creazione lista:} Spesso a casa (in cucina) ma anche sui mezzi pubblici o in università (quando si ha tempo morto e si cerca di fare mente locale). È qui che l'IA dà il massimo vantaggio, non potendo l'utente guardare il frigo.
    \item \textbf{Consultazione lista:} Esclusivamente al supermercato (spesso offline o con poca connessione, l'IA deve aver già generato la lista).
\end{itemize}

\subsection{Che relazione c'è tra utente e dati? (Privacy)}
L'app tratterà dati sulle abitudini alimentari. Dalle interviste (che farai) potrebbe emergere che agli utenti non importa condividere la frequenza con cui comprano la pasta, ma è fondamentale rassicurarli che i dati dell'IA sono usati localmente o in modo anonimo solo per migliorare le loro predizioni, e non venduti a supermercati per pubblicità mirata.

\subsection{Quali altri strumenti ha l'utente per completare il task?}
Esistono app come \textit{Bring!} o \textit{Any.do}, ma si basano su inserimento manuale. Esistono assistenti vocali ("Alexa, aggiungi latte alla lista"), ma richiedono un'azione proattiva dell'utente nel momento esatto in cui il prodotto finisce. Nessuno strumento attuale \textbf{anticipa} il bisogno basandosi sul consumo storico.

\subsection{Come comunicano gli utenti tra loro relativamente ai task?}
Gli utenti comunicano in modo asincrono. Un coinquilino finisce l'olio e (se si ricorda) lo scrive nel gruppo WhatsApp. L'IA di FarFromHome sostituirà questa comunicazione: se l'IA sa che una bottiglia d'olio dura in media 3 settimane per 3 persone, lo suggerirà automaticamente senza che i coinquilini debbano ricordarsi di comunicarlo.

\subsection{Con quale frequenza sono eseguiti i task?}
La generazione e il completamento della lista avvengono tipicamente 1 o 2 volte a settimana per la spesa principale, e quotidianamente per piccoli acquisti (es. pane). L'IA dovrà adattarsi a questi cicli.

\subsection{Quali sono i vincoli di tempo sui task, se ce ne sono?}
La revisione dei suggerimenti dell'IA deve richiedere meno di 30 secondi. Se l'utente deve perdere tempo a cancellare troppe previsioni errate dell'IA, l'app perde di utilità rispetto al foglio di carta.

\subsection{Che accade quando le cose vanno male durante l'esecuzione dei task?}
Se l'IA suggerisce prodotti sbagliati (falsi positivi) o non suggerisce qualcosa che manca (falsi negativi), l'utente interviene manualmente. Il sistema deve prevedere un meccanismo di \textit{feedback implicito} (se l'utente scarta un suggerimento, l'IA impara che ha sbagliato i calcoli e aggiusta i pesi per la prossima settimana).

\section{Annotazioni}
\textit{In questa sezione inserirai i veri fatti che hai osservato durante le interviste.}

\textbf{Esempio di Fatti osservati (da adattare con i tuoi veri risultati):}
\begin{itemize}
    \item Le persone trovano noioso dover controllare la scadenza o la quantità dei prodotti in dispensa.
    \item È emerso che i prodotti più dimenticati sono quelli di "back-up" (sale, olio, zucchero, carta igienica) perché non si consumano in un giorno, ma lentamente.
    \item Molti utenti single/fuorisede sono propensi a farsi guidare completamente da un'app se questo significa risparmiare 10 minuti di tempo.
    \item Le madri di famiglia intervistate si fidano meno dell'IA e preferiscono avere il controllo finale sulla lista, usando le predizioni solo come "promemoria" e non come lista definitiva.
    \item Se la precisione dell'IA è percepita come bassa nelle prime 2 settimane, l'utente abbandona la feature. Il \textit{cold-start} (i primi giorni di utilizzo dell'app) deve essere gestito bene, magari chiedendo all'inizio "Ogni quanto compri il latte?".
\end{itemize}

\end{document}
